\begin{longtable}[c]{|p{2.5cm}|p{13cm}|} 
    
  \caption{Segments of a Nominal Bit Time}
  \label{tab:twai-seg-of-a-nominal-bit-time} 
  \\\hline
  \rowcolor{lightgray}
  \textbf{Segment} & \textbf{Description}
  \\\hline
  \endfirsthead
  
  \multicolumn{2}{c}{\bfseries \tablename\ \thetable{} -- cont'd from previous page}\\\hline
 
  \rowcolor{lightgray}
    \textbf{Segment} & \textbf{Description}\\\hline
  \endhead
  \multicolumn{2}{r}{\textbf{Cont'd on next page}} \\
  \endfoot
  \hline
  \endlastfoot

    SS & The SS (Synchronization Segment) is 1 Time Quantum long. If all nodes are perfectly synchronized, the edge of a bit will lie in the SS.
            \\\hline
    PBS1 & PBS1 (Phase Buffer Segment 1) can be 1 to 16 Time Quanta long. PBS1 is meant to compensate for the physical delay times within the network. PBS1 can also be lengthened for synchronization purposes.
            \\\hline
    PBS2 & PBS2 (Phase Buffer Segment 2) can be 1 to 8 Time Quanta long. PBS2 is meant to compensate for the information processing time of nodes. PBS2 can also be shortened for synchronization purposes.
            \\\hline
    
\end{longtable}
